\begin{figure}[htbp]
	\centering
	\begin{tikzpicture}[>=Stealth, line width=0.8pt]
		
		% ═══════════════════════════════════════════════════════════
		% PANNELLO (a) — Prestazione cinematica del glifo esterno
		% ═══════════════════════════════════════════════════════════
		\begin{scope}[xshift=0cm]
			
			%% ── Parametri ──
			\def\xA{0.00}
			\def\yA{0.00}
			\def\Lpiatto{0.60}    % semi-lunghezza piatto
			\def\bw{0.20}         % semi-distanza ruote
			\def\rrota{0.05}      % raggio ruote
			\def\gw{0.65}         % semi-larghezza barra di terra (un po' più del piatto)
			\def\rrota{0.10}     % raggio ruote (tangenti ai due piatti)
			
			%% ── Asse a (retta verticale grigia continua) ──
			\draw[gray!60, line width=0.6pt] 
			(\xA, -0.55) -- (\xA, 2.85);
			
			%% ── Etichetta del CR all'infinito ──
			\node[gray, font=\footnotesize, anchor=south west] 
			at (\xA+0.30, 2.20) {$C = \infty_\ab$};
			
			%% ── Versore a in alto ──
			\draw[->, line width=1pt] 	(\xA, 2.10) -- (\xA, 2.85);
			\node[anchor=west, font=\small] 
			at (\xA+0.05, 2.85) {$\ab$};
			
			%% ── Corpo C (con piatto orizzontale inferiore) ──
			\draw[line width=1pt] 
			(-\Lpiatto, 0.00) 
			.. controls (-1.55, 0.10) and (-1.80, 0.95) 
			.. (-1.70, 1.55)
			.. controls (-0.85, 2.40) and ( 0.35, 2.50) 
			.. ( 0.85, 2.15)
			.. controls ( 1.65, 1.70) and ( 1.80, 0.60) 
			.. (\Lpiatto, 0.00)
			-- cycle;
			\node at (0.50, 1.60) {$\mathcal{C}$};
			
			%% ── Equazioni di vincolo (in alto, accanto all'asse) ──
			\node[anchor=west, font=\small] at (\xA-2.40, 2.55) {$\ub(A)\cdot\ab = 0$};
			\node[anchor=west, font=\small] at (\xA-2.40, 2.10) {$\vartheta = 0$};
			
			
			%% ── Glifo: prima barra (sotto il piatto del corpo) ──
			%\draw (-\gw, -0.05) -- (\gw, -0.05);
			
			%% ── Glifo: due ruote ──
			\draw (-\bw, -0.13) circle (\rrota);
			\draw (\bw, -0.13) circle (\rrota);
			
			%% ── Glifo: due ruote (tangenti ai due piatti) ──
			\draw (-\bw, -0.155) circle (\rrota);
			\draw (\bw, -0.155) circle (\rrota);
			
			%% ── Glifo: seconda barra (al telaio) ──
			\draw (-\gw, -0.26) -- (\gw, -0.26);
			
			%% ── Glifo: tratteggio del suolo ──
			\foreach \x in {-0.50,-0.25,0,0.25,0.50}
			\draw[thin] (\x, -0.26) -- ++(-0.10, -0.13);
			
			%% ── Punto A e label ──
			\fill (\xA, \yA) circle (1.5pt);
			\node[anchor=south west, font=\small] at (\xA, \yA+0.05) {$A$};
			
			
			%% ── Etichetta verticale "asse del glifo" ──
			\node[gray, rotate=90, anchor=south, font=\footnotesize] at (\xA, 1.20) {asse del glifo};
			
			
			%% ── Etichetta pannello ──
			\node[anchor=north, font=\small] 	at (\xA-1.95, -0.25) {$(a)$};
			
		\end{scope}
		
		% ═══════════════════════════════════════════════════════════
		% PANNELLO (b) — Prestazione statica del glifo esterno
		% ═══════════════════════════════════════════════════════════
		\begin{scope}[xshift=5cm]
			
			%% ── Parametri (identici ad (a)) ──
			\def\xA{0.00}
			\def\yA{0.00}
			\def\Lpiatto{0.60}
			\def\bw{0.20}
			\def\rrota{0.10}
			\def\gw{0.65}
			
			%% ── Asse a (retta verticale grigia continua) ──
			\draw[gray!60, line width=0.6pt] 
			(\xA, -0.55) -- (\xA, 2.85);
			
			%% ── Versore a in alto ──
			\draw[->, line width=1pt] 
			(\xA, 2.10) -- (\xA, 2.85);
			\node[anchor=west, font=\small] 
			at (\xA+0.05, 2.85) {$\ab$};
			
			%% ── Corpo C (identico ad (a)) ──
			\draw[line width=1pt] 
			(-\Lpiatto, 0.00) 
			.. controls (-1.55, 0.10) and (-1.80, 0.95) 
			.. (-1.70, 1.55)
			.. controls (-0.85, 2.40) and ( 0.35, 2.50) 
			.. ( 0.85, 2.15)
			.. controls ( 1.65, 1.70) and ( 1.80, 0.60) 
			.. (\Lpiatto, 0.00)
			-- cycle;
			\node at (0.50, 1.60) {$\mathcal{C}$};
			
			%% ── Glifo: due ruote tangenti ──
			\draw (-\bw, -0.155) circle (\rrota);
			\draw (\bw, -0.155) circle (\rrota);
			
			%% ── Glifo: seconda barra (al telaio) ──
			\draw (-\gw, -0.26) -- (\gw, -0.26);
			
			%% ── Glifo: tratteggio del suolo ──
			\foreach \x in {-0.50,-0.25,0,0.25,0.50}
			\draw[thin] (\x, -0.26) -- ++(-0.10, -0.13);
			
			%% ── Punto A e label ──
			\fill (\xA, \yA) circle (1.5pt);
			\node[anchor=south west, font=\small] 	at (\xA+0.1, \yA-0.04) {$A$};
			
			%% ──────────────────────────────────────────────────
			%% REAZIONE: forza R(A) a + coppia μ
			%% ──────────────────────────────────────────────────
			
			%% Forza R(A) a: freccia verticale lungo l'asse, da A verso l'alto
			\draw[->, line width=1.4pt] 
			(\xA, \yA+0.05) -- (\xA, \yA+1.20);
			\node[anchor=west, font=\small] at (\xA+0.08, \yA+1.0) {$R(A)\,\ab$};
			
			%% Coppia reattiva μ(A): semicerchio centrato in A, sopra il punto
			\def\rmu{0.30}
			\draw[->, line width=1.4pt] ({\xA+\rmu}, {\yA+0.30}) arc(0:180:\rmu);
			\node[anchor=west, font=\small] at (\xA-1.30, {\yA+\rmu+0.5}) {$\mu(A)$};
			
			%% ── Etichetta pannello ──
			\node[anchor=north, font=\small] at (\xA-1.95, -0.25) {$(b)$};
			
		\end{scope}
		
	\end{tikzpicture}
	\caption[Glifo esterno con asse $\ab$: prestazioni cinematica e statica.]{%
		Glifo esterno applicato al corpo $\mathcal{C}$ in $A$,
		con asse $\ab$. Prestazione cinematica (a): le equazioni
		di vincolo $\ub(A)\cdot\ab = 0$ e $\vartheta = 0$
		impediscono la traslazione di $A$ lungo $\ab$ e la
		rotazione del corpo; il \CRa\ è all'infinito sulla retta
		dell'asse, $C = \infty_\ab$. Prestazione statica (b):
		la reazione vincolare consta di una forza $R(A)\,\ab$
		diretta lungo l'asse e di una coppia $\mu(A)$, in
		dualità rispettivamente con la traslazione lungo $\ab$
		e con la rotazione bloccate, in accordo con la
		molteplicità due.}
	\label{fig:glifo_esterno}
\end{figure}
